% ============================================================
% Capítulo 0 — Prólogo
% ============================================================
\chapter{Prólogo}

Este libro es una guía exhaustiva del proyecto \texttt{textreid-train}:
un repositorio limpio, enfocado exclusivamente al entrenamiento y la
evaluación del modelo \textbf{TextReIDNet} para \emph{reidentificación
de personas basada en texto}. No es un manual de uso ni un
\emph{quickstart}; es un documento de referencia técnica que recorre
cada archivo del proyecto y explica, capa por capa, qué hace y por qué.

\section{¿A quién está dirigido?}

\begin{itemize}
  \item \textbf{Estudiantes o practitioners de visión por computador}
        que quieran entender un sistema real de
        \emph{re-ID} basado en texto, desde la matemática del paper
        hasta los detalles de implementación.
  \item \textbf{Investigadores} que busquen reproducir, comparar o
        extender el modelo.
  \item \textbf{Ingenieros} que tengan que mantener, depurar o
        desplegar este código en producción.
\end{itemize}

\section{Qué asume el lector}

\begin{itemize}
  \item Familiaridad básica con \textbf{PyTorch} (tensores,
        \texttt{nn.Module}, \texttt{DataLoader}).
  \item Conocimiento general de \emph{deep learning}: redes
        convolucionales, recurrentes, funciones de pérdida y
        optimización.
  \item Conocimiento elemental de procesamiento de lenguaje natural:
        \emph{tokenización}, \emph{embeddings}, \acronym{GRU}{GRU}.
\end{itemize}

No se asume experiencia previa con \emph{reidentificación de personas},
\emph{EfficientNet}, \emph{BERT} ni métricas como \acronym{mAP}{mAP}.

\section{Convenciones tipográficas}

\begin{itemize}
  \item Los nombres de archivos se escriben en
        \file{\texttt{monoespaciado}} y se citan con su ruta relativa
        al raíz del proyecto: \file{model/textreidnet.py}.
  \item Las referencias a líneas específicas se hacen con el formato
        \fline{archivo.py}{L}. Por ejemplo, \fline{config.py}{42}
        apunta a la línea 42 de \file{config.py}.
  \item Las formas de los tensores se muestran entre paréntesis en
        color azul: \dims{(B, 1280, 16, 16)} significa
        \emph{batch} × \emph{canales} × \emph{alto} × \emph{ancho}.
  \item Los parámetros de configuración aparecen en
        \texttt{\textbf{negrita mono}}.
  \item Los listados de código van numerados y muestran sólo la región
        relevante. La numeración coincide con la del archivo
        original.
\end{itemize}

\section{Estructura del libro}

El libro se divide en cuatro partes:

\begin{enumerate}
  \item \textbf{Fundamentos teóricos} (Caps.~\ref{cap:introduccion}
        y~\ref{cap:arquitectura-teorica}): contexto del problema,
        formulación matemática y arquitectura de U-TextReIDNet según
        el paper.
  \item \textbf{Sistema descentralizado} (Cap.~\ref{cap:sistema-descentralizado}):
        cómo se despliega el modelo en una red de cámaras reales.
  \item \textbf{Implementación} (Caps.~\ref{cap:dataset}--\ref{cap:evaluacion}):
        dataset, configuración, redes, pérdidas, entrenamiento y
        evaluación. Es el grueso del libro.
  \item \textbf{Apéndices} (Cap.~\ref{cap:apendices}): glosario,
        tablas de referencia rápida.
\end{enumerate}

\section{Código fuente}

Todos los listados del libro provienen de los archivos reales del
repositorio \texttt{textreid-train}. Se recomienda abrir el repositorio
en paralelo a la lectura para poder navegar entre el PDF y el código
sin perderse.

\begin{center}
\begin{tikzpicture}
  \node[draw, rounded corners, fill=codebg, minimum width=6cm, align=center]
    {\textbf{\large textreid-train/}\\[2pt]
     \texttt{model/} · \texttt{datasets/} · \texttt{evaluation/}\\[1pt]
     \texttt{utils/} · \texttt{scripts/} · \texttt{config.py}};
\end{tikzpicture}
\end{center}

\section*{Sobre este documento}

\noindent\textbf{Versión}: 1.0\\
\textbf{Fecha de compilación}: \today\\
\textbf{Paper de referencia}: Agyeman \emph{et al.}, ``Decentralized
Text-Based Person Re-Identification in Multicamera Networks'',
\emph{IEEE Access}, noviembre de 2024. DOI:~\href{https://doi.org/10.1109/ACCESS.2024.3501382}{10.1109/ACCESS.2024.3501382}.\\
\textbf{Traducción del paper al español}:
\file{docs/paper.md} (552 líneas).

\vspace{1em}
\noindent\rule{\textwidth}{0.4pt}
\vspace{0.5em}

\noindent\textit{Conviene leer este libro con un navegador de código
abierto. Cada vez que veas una referencia como
\fline{model/textreidnet.py}{13}, podrás saltar directamente al
constructor de la clase \texttt{TextReIDNet}.}
